% !TEX root = ./main.tex

\DeclareMathAlphabet{\mathpzc}{OT1}{pzc}{m}{it}

%%%%    TODONOTES    %%%%
\newcommand{\todoin}[1]{\todo[inline]{#1}}
\newcommand{\todoinline}[1]{\todoin{#1}}

\newcommand{\genericcomment}[3]{\todo[color=#1,size=\tiny,fancyline,author=#2]{#3}\xspace}
\newcommand{\genericcommentin}[3]{\todo[inline,color=#1,size=\footnotesize,author=#2]{#3}\xspace}
\newcommand{\genericatsomebody}[2]{\colorbox{#1}{\textbf{@#2}}}

\newcommand{\lv}[1]{\genericcomment{violet!35}{Lena}{#1}\xspace}
\newcommand{\lvin}[1]{\genericcommentin{violet!35}{Lena}{#1}\xspace}
\newcommand{\atlv}{\genericatsomebody{violet!35}{Lena}\xspace}

\newcommand{\tg}[1]{\todo[color=cyan!30,size=\tiny,fancyline,author=Tobias]{#1}\xspace}
\newcommand{\tgin}[1]{\todo[inline,color=cyan!30,size=\footnotesize,author=Tobias]{#1}\xspace}
\newcommand{\attg}{\colorbox{cyan!30}{\textbf{@Tobias}}}

\newcommand{\blk}[1]{\todo[color=orange!60,size=\tiny,fancyline,author=Benjamin]{#1}\xspace}
\newcommand{\blkin}[1]{\todo[inline,color=orange!60,size=\footnotesize,author=Benjamin]{#1}\xspace}
\newcommand{\atblk}{\colorbox{orange!60}{\textbf{@Benjamin}}}



%%%% FONTS %%%%

\newcommand{\sfsymbol}[1]{\textsf{\upshape {#1}}}
\newcommand{\ttsymbol}[1]{\texttt{\upshape {#1}}}



%%%% PBLIMP Specifics %%%%
\newcommand{\Varso}{\sfsymbol{Vars}_o}
\newcommand{\Varsu}{\sfsymbol{Vars}_u}



%%%% Standards %%%%

\newcommand{\sk}{\mathfrak{s}} % stack
\newcommand{\hh}{\mathfrak{h}} % heap
\newcommand{\skhh}{(\sk,\hh)} % generic state
%\newcommand{\Stacks}{\mathcal{S}}
\newcommand{\Stacks}{\mathsf{Stacks}}
%\newcommand{\Heaps}{\mathcal{H}}
\newcommand{\Heaps}{\mathsf{Heaps}}
\newcommand{\emptyheap}{\hh_\emptyset}
\newcommand{\singleheap}[2]{\{ #1 \,\mapsto\, #2\}}
\newcommand{\loc}{\ell}
\newcommand{\val}{\vala}
\newcommand{\vala}{v}
\newcommand{\valb}{w}
\newcommand{\Loc}{L}
\newcommand{\disjoint}{\mathrel{\bot}}
%\newcommand{\States}{\Sigma}
\newcommand{\States}{\mathsf{States}}
\newcommand{\To}{\rightarrow}
\newcommand{\true}{\mathsf{true}}
\newcommand{\false}{\mathsf{false}}
\newcommand{\mydot}{\text{{\Large\textbf{.}}~}}
\newcommand{\guard}{\varphi}
\newcommand{\heapSize}{\sfsymbol{\textbf{size}}}

%%%% Operational Semantics %%%
\newcommand{\transitionRelation}{\mathrel{{\large \rightsquigarrow}}}
\newcommand{\transition}[2]{#1  \transitionRelation  #2 }
\newcommand{\noTransition}[2]{#1 \not\transitionRelation #2}
\newcommand{\transitionRelationClosure}{\transitionRelation^{*}}
\newcommand{\configuration}[2]{\left\langle #1, #2 \right\rangle}
\newcommand{\toState}[2]{#1 ~ {}{\large \transitionRelationClosure}{} ~ #2}
\newcommand{\notToState}[2]{#1 \mathrel{{\large \not\transitionRelation}^{*}} #2}
\newcommand{\nonTermination}{\uparrow}
\newcommand{\abort}{\textbf{abort}}
\newcommand{\evaluateExpr}[2]{#2(#1)}
\newcommand{\updateHeap}[2]{\left[ h \mid #2 \right]}
\newcommand{\terminated}{\ensuremath{\downarrow}}
\newcommand{\terminatedState}[1]{\configuration{\terminated}{#1}}


%%%%    Shortcuts    %%%%

\newcommand{\qif}{\quad\textnormal{if}~}
\newcommand{\tif}{~\textnormal{if}~}

\newcommand{\qiff}{\quad\textnormal{iff}\quad}
\newcommand{\qqiff}{\qquad\textnormal{iff}\qquad}
\newcommand{\lqiff}{\textnormal{iff}\quad}
\newcommand{\lqqiff}{\textnormal{iff}\qquad}

\newcommand{\qand}{\quad\textnormal{and}\quad}
\newcommand{\qqand}{\qquad\textnormal{and}\qquad}

\newcommand{\qor}{\quad\textnormal{or}\quad}
\newcommand{\qqor}{\qquad\textnormal{or}\qquad}

\newcommand{\qimplies}{\quad\textnormal{implies}\quad}
\newcommand{\qqimplies}{\qquad\textnormal{implies}\qquad}
\newcommand{\lqimplies}{\textnormal{implies}\quad}
\newcommand{\lqqimplies}{\textnormal{implies}\qquad}

\newcommand{\pprec}{\mathrel{{\prec}\hspace{-.5ex}{\prec}}}
\newcommand{\ppprec}{~{}\pprec{}~}


\newcommand{\morespace}[1]{~{}#1{}~}
\newcommand{\qmorespace}[1]{\quad{}#1{}\quad}
\newcommand{\qqmorespace}[1]{\qquad{}#1{}\qquad}
\newcommand{\ppreceq}{~{}\preceq{}~}
\newcommand{\ssucceq}{~{}\succeq{}~}
\newcommand{\eeq}{~{}={}~}
\newcommand{\ttriangleq}{~{}\triangleq{}~}
\newcommand{\defeq}{\coloneqq}
\newcommand{\ddefeq}{~{}\coloneqq{}~}
\newcommand{\neeq}{~{}\neq{}~}
\newcommand{\qeq}{\quad{}={}\quad}
\newcommand{\qqeq}{\qquad{}={}\qquad}
\newcommand{\lleq}{~{}\leq{}~}
\newcommand{\LL}{~{}<{}~}
\newcommand{\GG}{~{}>{}~}
\newcommand{\ggeq}{~{}\geq{}~}
\newcommand{\ssqsubseteq}{~{}\sqsubseteq{}~}
\newcommand{\iimplies}{~{}\implies{}~}
\newcommand{\mmodels}{~{}\models{}~}
\newcommand{\nnotmodels}{~{}\not\models{}~}
\newcommand{\iin}{~{}\in{}~}
\newcommand{\nnotin}{~{}\not\in{}~}
\newcommand{\ttimes}{\morespace{\times}}
\newcommand{\tto}{\morespace{\to}}

\newcommand{\ttriangleright}{\morespace{\triangleright}}
\newcommand{\qtriangleright}{\qmorespace{\triangleright}}
\newcommand{\qqtriangleright}{\qqmorespace{\triangleright}}

\newcommand{\qedtriangle}{\hfill $\triangle$}

\newcommand{\RSL}{\sfsymbol{RSL}\xspace}
\newcommand{\SL}{\sfsymbol{SL}\xspace}

%%%%    SYMBOLS    %%%%
\newcommand{\cmark}{{\large\ding{51}}}
\newcommand{\xmark}{}%{\scriptsize\ding{55}}}

%%%%    LOGICS    %%%%
\newcommand{\slemp}{\sfsymbol{\textbf{emp}}}
\newcommand{\emp}{\iverson{\slemp}}
\newcommand{\iemp}{\iiverson{\slemp}}
%\newcommand{\tree}[1]{\iverson{\sfsymbol{\textbf{\text{tree}(\(\boldsymbol{#1}\))}}}}
%\newcommand{\bst}[1]{\iverson{\sfsymbol{\textbf{\text{bst}(\(\boldsymbol{#1}\))}}}}
%\newcommand{\elem}[1]{\iverson{\sfsymbol{\textbf{\text{elem}(\(\boldsymbol{#1}\))}}}}
%\newcommand{\contains}[1]{\iverson{\sfsymbol{\textbf{\text{contains}(\(\boldsymbol{#1}\))}}}}
%\newcommand{\dll}[2]{\iverson{\sfsymbol{\textbf{\text{dll}(\(\boldsymbol{#1, #2}\))}}}}
%\newcommand{\rroot}[2]{\iverson{\sfsymbol{\textbf{\text{root}(\(\boldsymbol{#1, #2}\))}}}}
%\newcommand{\suc}[2]{\iverson{\sfsymbol{\textbf{\text{succ}(\(\boldsymbol{#1, #2}\))}}}}
\newcommand{\slsingleton}[2]{{#1 \mapsto #2}}
\newcommand{\slvalidpointer}[1]{{#1 \mapsto \,{-}\,}}

\newcommand{\singleton}[2]{\iverson{#1 \mapsto #2}}
\newcommand{\containsPointer}[2]{\iverson{#1 \hookrightarrow #2}}
\newcommand{\notContainsPointer}[2]{\iverson{#1 \not\hookrightarrow #2}}
\newcommand{\validpointer}[1]{\iverson{#1 \mapsto \,{-}\,}}

\newcommand{\isingleton}[2]{\iiverson{#1 \mapsto #2}}
\newcommand{\icontainsPointer}[2]{\iiverson{#1 \hookrightarrow #2}}
\newcommand{\inotContainsPointer}[2]{\iiverson{#1 \not\hookrightarrow #2}}
\newcommand{\ivalidpointer}[1]{\iiverson{#1 \mapsto \,{-}\,}}

\newcommand{\sepcon}{\mathbin{{\star}}} 
\newcommand{\sepadd}{\mathbin{{\oplus}}}
\newcommand{\ssepadd}{\morespace{\sepadd}}
\newcommand{\sepimp}{\mathbin{\text{\raisebox{-0.1ex}{$\boldsymbol{{-}\hspace{-.55ex}{-}}$}}\hspace{-1ex}\text{\raisebox{0.13ex}{\rotatebox{-17}{$\star$}}}}}
\newcommand{\sepmon}{\mathbin{\text{\raisebox{-0.0ex}{$\boldsymbol{{-}\hspace{-.55ex}{-}}$}}\hspace{-0.3ex}\text{$\ominus$}}}
\newcommand{\ssepmon}{\morespace{\sepmon}}
\newcommand{\nil}{\textit{nil}}
\newcommand{\nt}[1]{\textit{numtree}\left(#1\right)}

\newcommand{\emprun}[1]{\ensuremath{\underline{#1}}}
%\newcommand{\pureemp}[1]{\ensuremath{\left\langle #1 \wedge \slemp \right\rangle}}
\newcommand{\pureemp}[1]{\emprun{\iiverson{#1}}}

%%%% Potential function %%%%

\newcommand{\htree}[1]{\iverson{\sfsymbol{\textbf{\(\hat{\text{tree}}(\boldsymbol{#1})\)}}}}
\newcommand{\hsuc}[2]{\iverson{\sfsymbol{\textbf{\(\hat{\text{succ}}(\boldsymbol{#1, #2})\)}}}}
\newcommand{\pii}[1]{\sfsymbol{\(\boldsymbol{\vartheta(#1)}\)}}
\newcommand{\sigmaa}[1]{\sfsymbol{\(\boldsymbol{\sigma(#1)}\)}}
\newcommand{\hpii}[1]{\sfsymbol{\(\boldsymbol{\hat{\vartheta}(#1)}\)}}
\newcommand{\hsigmaa}[1]{\sfsymbol{\(\boldsymbol{\hat{\sigma}(#1)}\)}}



%%%%    WP-STYLE TRANSFORMERS    %%%%

\newcommand{\hoare}[3]{\left\langle \,{#1}\vphantom{#3}\, \right\rangle \mathrel{#2} \left\langle \, {#3}\vphantom{#1} \, \right\rangle}
\newcommand{\partialhoare}[3]{\left\langle \,{#1}\vphantom{#3}\, \right\rangle \mathrel{\widetilde{#2}} \left\langle \, {#3}\vphantom{#1} \, \right\rangle}

% general transformer
\newcommand{\transformersymbol}{\mathcal{T}}
\newcommand{\boldtransformersymbol}{\boldsymbol{\transformersymbol}}
\newcommand{\transformer}[2]{\transformersymbol \, \left\llbracket {#1} \right\rrbracket \, \left( {#2} \right)}
\newcommand{\transformerC}[1]{\transformersymbol \, \left\llbracket {#1} \right\rrbracket}

% weakest preexpectation
\newcommand{\wpsymbol}{\sfsymbol{wp}}
\newcommand{\boldwpsymbol}{\textbf{\sfsymbol{wp}}}
\let\weierwp\wp
\renewcommand{\wp}[2]{\wpsymbol \left\llbracket {#1} \right\rrbracket \left( {#2} \right)}
\newcommand{\wpC}[1]{\wpsymbol \, \left\llbracket {#1} \right\rrbracket}
\newcommand{\cemphwp}[2]{\cemph{\wpsymbol} \, \left\llbracket {#1} \right\rrbracket \, \left( {#2} \right)}

\newcommand{\wpk}[3]{\wpsymbol^{= #3} \, \left\llbracket {#1} \right\rrbracket \, \left( {#2} \right)}

% angelic weakest preexpectation
\newcommand{\awpsymbol}{\sfsymbol{awp}}
\newcommand{\boldawpsymbol}{\textbf{\sfsymbol{awp}}}
\newcommand{\awp}[2]{\awpsymbol \, \left\llbracket {#1} \right\rrbracket \, \left( {#2} \right)}
\newcommand{\awpC}[1]{\awpsymbol \, \left\llbracket {#1} \right\rrbracket}
\newcommand{\cemphawp}[2]{\cemph{\awpsymbol} \, \left\llbracket {#1} \right\rrbracket \, \left( {#2} \right)}

% weakest liberal preexpectation
\newcommand{\wlpsymbol}{\sfsymbol{wlp}}
\newcommand{\boldwlpsymbol}{\textbf{\sfsymbol{wlp}}}
\newcommand{\wlp}[2]{\wlpsymbol \, \left\llbracket {#1} \right\rrbracket \, \left( {#2} \right)}
\newcommand{\cemphwlp}[2]{\cemph{\wlpsymbol} \, \left\llbracket {#1} \right\rrbracket \, \left( {#2} \right)}
\newcommand{\wlpC}[1]{\wlpsymbol \, \left\llbracket {#1} \right\rrbracket}

% angelic weakest liberal preexpectation
\newcommand{\awlpsymbol}{\sfsymbol{awlp}}
\newcommand{\boldawlpsymbol}{\textbf{\sfsymbol{awlp}}}
\newcommand{\awlp}[2]{\awlpsymbol \, \left\llbracket {#1} \right\rrbracket \, \left( {#2} \right)}
\newcommand{\awlpC}[1]{\awlpsymbol \, \left\llbracket {#1} \right\rrbracket}
\newcommand{\cemphawlp}[2]{\cemph{\awlpsymbol} \, \left\llbracket {#1} \right\rrbracket \, \left( {#2} \right)}

% expected runtime transformer
\newcommand{\ertsymbol}{\sfsymbol{ert}}
\newcommand{\boldertsymbol}{\textbf{\sfsymbol{ert}}}
\newcommand{\ert}[2]{\ertsymbol \, \left\llbracket {#1} \right\rrbracket \, \left( {#2} \right)}
\newcommand{\ertC}[1]{\ertsymbol \, \left\llbracket {#1} \right\rrbracket}
\newcommand{\ertk}[3]{\ertsymbol^{#3} \, \left\llbracket {#1} \right\rrbracket \, \left( {#2} \right)}

% conditional weakest preexpectation
\newcommand{\cwpsymbol}{\sfsymbol{cwp}}
\newcommand{\qcwpsymbol}{\sfsymbol{\underline{cw}p}}

%\newcommand{\conditionalPair}[2]{{\let\oldarraystretch\arraystretch}\renewcommand{\arraystretch}{1}~\holter{~\raisebox{0ex}{${#1}$}~}{~\raisebox{0.25ex}{${#2}$}~}~\renewcommand{\arraystretch}{\oldarraystretch}}

\newcommand{\conditionalPair}[2]{\inlineConditionalPair{#1}{#2}}
\newcommand{\smallconditionalPair}[2]{\inlineConditionalPair{\text{\scriptsize $#1$}}{\text{\scriptsize $#2$}}}

\newcommand{\tikzholtpoltconditionalPair}[2]{%
    \tikz[baseline=(center)]{
\node[inner xsep=.3em, inner ysep=.3em, outer sep=0em] (todotted) {$\displaystyle{#1}$};
\node[inner xsep=.3em, inner ysep=.3em, outer sep=0em, below=0em of todotted, yshift=-0em] (todotted2) {$\displaystyle#2$};
\node (upperbox) [inner sep=0em, outer sep=0em, draw=none, fit = (todotted.north west) (todotted2.north west) (todotted.south east) (todotted2.north east)] {};
\node (lowerbox) [inner sep=0em, outer sep=0em, draw=none, fit = (todotted.south west) (todotted2.north west) (todotted.south east) (todotted2.south east)] {};
\draw[use as bounding box, draw=none] ($(upperbox.north west) + (-.1em, -.2em)$) rectangle ($(lowerbox.south east) + (.1em, .2em)$);
\node (center) at ($(upperbox.south) + (0, -0.5ex)$) {};
\draw[semithick] (upperbox.north west) -- (upperbox.south west) -- (upperbox.south east) -- (lowerbox.south east);
    }%
}%


\newcommand{\inlineConditionalPair}[2]{%
    \tikz[baseline=(center)]{
\node[inner xsep=.1em, inner ysep=.2em, outer sep=0em, draw=none, auto] (todotted) {\small $#1\!\!$};
\node[inner xsep=.1em, inner ysep=.2em, outer sep=0em, right=.5em of todotted, yshift=-.3em, draw=none, auto] (todotted2) {\small $#2$};
\draw[use as bounding box, draw=none] ($(todotted.north west) + (0, -.3em)$) rectangle ($(todotted2.south east) + (0, .3em)$);
\node (center) at ($(todotted.base)!0.5!(todotted2.base)$) {};
\draw (todotted.south west) -- (todotted.south east) -- (todotted2.north west) -- (todotted2.north east);
    }%
}%

\newcommand{\boldcwpsymbol}{\textbf{\sfsymbol{cwp}}}
\newcommand{\cwponearg}[2]{\cwpsymbol \, \left\llbracket {#1} \right\rrbracket \, \left( {#2} \right)}
\newcommand{\cwp}[3]{\cwponearg{#1}{\inlineConditionalPair{#2}{#3}}}
\newcommand{\cwpC}[1]{\cwpsymbol \, \left\llbracket {#1} \right\rrbracket}

\newcommand{\qcwp}[2]{\qcwpsymbol \, \left\llbracket {#1} \right\rrbracket \, \left( {#2} \right)}
\newcommand{\qcwpC}[1]{\qcwpsymbol \, \left\llbracket {#1} \right\rrbracket}

% conditional weakest liberal preexpectation
\newcommand{\cwlpsymbol}{\sfsymbol{cwlp}}
\newcommand{\qcwlpsymbol}{\sfsymbol{\underline{cwl}p}}
\newcommand{\boldcwlpsymbol}{\textbf{\sfsymbol{cwlp}}}
\newcommand{\cwlponearg}[2]{\cwlpsymbol \, \left\llbracket {#1} \right\rrbracket \, \left( {#2} \right)}
\newcommand{\cwlp}[3]{\cwlponearg{#1}{\inlineConditionalPair{#2}{#3}}}
\newcommand{\cwlpC}[1]{\cwlpsymbol \, \left\llbracket {#1} \right\rrbracket}

\newcommand{\qcwlp}[2]{\qcwlpsymbol \, \left\llbracket {#1} \right\rrbracket \, \left( {#2} \right)}
\newcommand{\qcwlpC}[1]{\qcwlpsymbol \, \left\llbracket {#1} \right\rrbracket}


% mixed-sign wp

\newcommand{\littlebigLbag}{\raisebox{-.1em}{\text{\Large $\Lbag$}}}
\newcommand{\littlebigRbag}{\raisebox{-.1em}{\text{\text{\Large $\Rbag$}}}}
\newcommand{\bigLbag}{\raisebox{-.2em}{\text{\LARGE $\Lbag$}}}
\newcommand{\bigRbag}{\raisebox{-.2em}{\text{\text{\LARGE $\Rbag$}}}}
\newcommand{\BigLbag}{\raisebox{-.2em}{\text{\huge $\Lbag$}}}
\newcommand{\BigRbag}{\raisebox{-.2em}{\text{\text{\huge $\Rbag$}}}}
\newcommand{\hugeLbag}{\raisebox{-.4em}{\text{\Huge $\Lbag$}}}
\newcommand{\hugeRbag}{\raisebox{-.4em}{\text{\text{\Huge $\Rbag$}}}}

%\newcommand{\wpp}[3]{\ensuremath{\wpsymbol\left[{#1}\right]\iwpair{#2}{#3}}}
\newcommand{\wpp}[2]{\wpsymbol \, \left\llbracket{#1} \right\rrbracket \, \left(#2\right)}

\newcommand{\iwpsymbol}{\ensuremath{{\textnormal{\textsf{iwp}}}}}
\newcommand{\boldiwpsymbol}{\ensuremath{\boldsymbol{{\textnormal{\textbf{\textsf{iwp}}}}}}}
\newcommand{\iwponearg}[2]{\iwpsymbol \, \left\llbracket{#1} \right\rrbracket \, \left({#2}\right)}
\newcommand{\iwp}[3]{\iwpsymbol \, \left\llbracket{#1} \right\rrbracket \, \EqPair{#2}{#3}}
\newcommand{\iwpC}[1]{\iwpsymbol \, \left\llbracket {#1} \right\rrbracket}
\newcommand{\boldiwp}[3]{\boldsymbol{\boldiwpsymbol \, \left\llbracket{#1} \right\rrbracket \, \EqPair{#2}{#3}}}



%%%%    PGCL PROGRAMS    %%%%
\newcommand{\SKIP}{\ttsymbol{skip}}
\newcommand{\EMPTY}{\ensuremath{\textnormal{\texttt{empty}}}}
\newcommand{\HALT}{\ensuremath{\textnormal{\texttt{halt}}}}
\newcommand{\ABORT}{\ensuremath{\textnormal{\texttt{abort}}}}
\newcommand{\DIVERGE}{\ensuremath{\textnormal{\texttt{diverge}}}}
\newcommand{\TICKsymbol}{\ensuremath{\textnormal{\texttt{tick}}}}
\newcommand{\TICK}[1]{\ensuremath{\TICKsymbol\left({#1}\right)}}


%\newcommand{\OBSERVE}[1]{\ensuremath{\textnormal{\texttt{observe}} ~ {#1}}}
\newcommand{\ASSUME}[1]{\ensuremath{\textnormal{\texttt{assume}} \, \left({#1}\right)}}
\newcommand{\ASSERT}[1]{\ensuremath{\textnormal{\texttt{assert}} \, \left({#1}\right)}}
\newcommand{\CALL}[1]{\ensuremath{\textnormal{\texttt{call}} ~ {#1}}}

\newcommand{\AssignSymbol}{\coloneqq}
\newcommand{\ASSIGN}[2]{\ensuremath{#1 \AssignSymbol #2}}
\newcommand{\CHOOSE}[2]{\ASSIGN{#1}{\texttt{choose}\left(#2\right)}}
\newcommand{\BMODEL}[2]{\ensuremath{\mathit{Pr}_{#1}\left(#2\right)}}
\newcommand{\SMODEL}[2]{\ensuremath{#1\left(#2\right)}}
\newcommand{\UNIFASSIGN}[3]{\ensuremath{\ASSIGN{#1}{\ttsymbol{Unif} \left(#2,#3\right)}}}

\newcommand{\AppAssignSymbol}{\mathrel{\textnormal{\texttt{:}}\hspace{-.1em}{\approx}}}
\newcommand{\PASSIGN}[2]{\ensuremath{#1 \AppAssignSymbol\hspace{.1em} #2}}

% hpGCL
\newcommand{\ALLOC}[2]{\ensuremath{{#1} \AssignSymbol \mathtt{alloc}\left( #2 \right)}}
% \newcommand{\dereference}[1]{\relationrightproject{\,#1\,}}
\newcommand{\ADDRESS}[1]{\langle#1\rangle}
\newcommand{\dereference}[1]{\ADDRESS{#1}}
\newcommand{\HASSIGN}[2]{\ensuremath{\dereference{#1} \AssignSymbol #2}}
\newcommand{\ASSIGNH}[2]{\ensuremath{#1 \AssignSymbol \dereference{#2}}}
\newcommand{\HASSIGNH}[2]{\ensuremath{\dereference{#1} \AssignSymbol \dereference{#2}}}
\newcommand{\LOOKUP}[2]{\ASSIGNH{#1}{#2}}
\newcommand{\MUTATE}[2]{\HASSIGN{#1}{#2}}
\newcommand{\FREE}[1]{\ensuremath{\mathtt{free}(#1)}}


\newcommand{\COMPOSE}[2]{\ensuremath{{#1}{\,\fatsemi}~ {#2}}}
\newcommand{\PCHOICE}[3]{\ensuremath{\left\{\, {#1} \,\right\}\mathrel{\left[{#2}\right]}\left\{\, {#3} \,\right\}}}
\newcommand{\NDCHOICE}[2]{\ensuremath{\left\{\, {#1} \,\right\}\mathrel{\Box}\left\{\, {#2} \,\right\}}}
\newcommand{\NDBRACK}{\ensuremath{\left.\right\}\mathrel{\Box}\left\{\right.}}


\newcommand{\IFSYMBOL}{\ensuremath{\textnormal{\texttt{if}}}}
\newcommand{\IF}[1]{\ensuremath{\IFSYMBOL\,\left(\, {#1} \,\right)\,\{}}
\newcommand{\ELSESYMBOL}{\ensuremath{\textnormal{\texttt{else}}}}
\newcommand{\ELSE}{\ensuremath{\}\,\ELSESYMBOL\,\{}}
\newcommand{\ITE}[3]{\ensuremath{\IFSYMBOL\;{#1}\;\texttt{then}\; #2 \; \ELSESYMBOL\; #3  }}
%\newcommand{\ITE}[3]{\ensuremath{\textttt{if}, {#1} \texttt{then} {#2} \,\ELSESYMBOL {#3} \,\right\}}}
\newcommand{\INFER}[3]{\ensuremath{\texttt{infer}\,\left(\, {#1} \,\right)\,\left\{\, {#2} \,\right\}\,\ELSESYMBOL\,\left\{\, {#3} \,\right\}}}
\newcommand{\ELIF}[1]{\ensuremath{\}\, \ELSESYMBOL\, \IFSYMBOL\, \left(\, {#1} \,\right)\,\{}}


\newcommand{\observefail}{\text{\lightning}}
\newcommand{\observesymbol}{\textnormal{\texttt{observe}}}
\newcommand{\OBSERVE}[1]{\ensuremath{\observesymbol (#1)}}
\newcommand{\OBSERVEASSIGN}[2]{\ensuremath{#1 \coloneqq \observesymbol \,\left(#2\right)}}

\newcommand{\WHILESYMBOL}{\ensuremath{\textnormal{\texttt{while}}}}
\newcommand{\WHILE}[1]{\ensuremath{\WHILESYMBOL \, \left(\, {#1} \,\right)\left\{\right.}}

\newcommand{\WHILEDO}[2]{\ensuremath{\WHILESYMBOL \, \left(\, {#1} \,\right)\left\{\, {#2} \,\right\}}}

\newcommand{\PROCEDURE}{\ensuremath{\textnormal{\texttt{\textbf{proc}}}}~}

%\newcommand{\BOUNDEDWHILE}[3]{\ensuremath{\WHILE^{<{#1}}\:(#2)\:\{#3\}}}
\newcommand{\BOUNDEDWHILE}[3]{\ensuremath{\WHILESYMBOL^{{#1}}\:(#2)\:\{#3\}}}


%%% LOGIC NAMES %%%
\newcommand{\QSL}{\textnormal{\sfsymbol{QSL}}\xspace} 



%%%%    SETS    %%%%
\newcommand{\PBLIMP}{\sfsymbol{\mbox{P-BLIMP}}\xspace}   % Programs
\newcommand{\gcl}{\textnormal{\sfsymbol{GCL}}\xspace}   % Programs
\newcommand{\pgcl}{\textnormal{\sfsymbol{pGCL}}\xspace}   % Programs
\newcommand{\cpgcl}{\textnormal{\sfsymbol{cpGCL}}\xspace}   % Programs
\newcommand{\hpgcl}{\textnormal{\sfsymbol{hpGCL}}\xspace}   % Programs
\newcommand{\boldgcl}{\textnormal{\textbf{\sfsymbol{GCL}}}\xspace}   % Programs
\newcommand{\boldhpgcl}{\textnormal{\textbf{\sfsymbol{hpGCL}}}\xspace}   % Programs
\newcommand{\boldpgcl}{\textnormal{\textbf{\sfsymbol{pGCL}}}\xspace}   % Programs
\newcommand{\boldcpgcl}{\textnormal{\textbf{\sfsymbol{cpGCL}}}\xspace}   % Programs
\newcommand{\Vars}{\ensuremath{\mathsf{Vars}}\xspace}   % Variables
\newcommand{\Vals}{\ensuremath{\mathsf{Vals}}\xspace}    % Values
%\newcommand{\Vals}{\ensuremath{\Nats}\xspace}    % Values
\newcommand{\Locs}{\ensuremath{\mathsf{{Loc}}\xspace}}    % memory locations

\newcommand{\ppdl}{\textnormal{\sfsymbol{PPDL}}\xspace}


\newcommand{\Nats}{\ensuremath{\mathbb{N}}\xspace}
\newcommand{\Ints}{\ensuremath{\mathbb{Z}}\xspace}
\newcommand{\Rats}{\ensuremath{\mathbb{Q}}\xspace}
\newcommand{\Reals}{\ensuremath{\mathbb{R}}\xspace}
\newcommand{\PosReals}{\mathbb{R}_{\geq 0}}
\newcommand{\StrictPosReals}{\mathbb{R}_{> 0}}
\newcommand{\PosRealsInf}{\mathbb{R}_{\geq 0}^\infty}
\newcommand{\PosRats}{\mathbb{Q}_{\geq 0}}
\newcommand{\PosRatsInf}{\mathbb{Q}_{\geq 0}^\infty}

\newcommand{\Preds}{\pot{\States}}
%\newcommand{\States}{\Sigma}
%\newcommand{\Confs}{\mathbb{K}}
\newcommand{\Confs}{\textsf{Conf}}
\newcommand{\E}{\mathbb{E}}
\newcommand{\Ebounded}{\mathbb{E}_{\leq {\, \exists b}}}
\newcommand{\Eone}{\mathbb{E}_{\leq 1}}
\newcommand{\Epm}{\ensuremath{\mathbb{E}_{{\pm}}}}
\newcommand{\A}{\mathbb{A}}
\newcommand{\Api}[1][\potent]{\A_{#1}}
\newcommand{\T}{\mathbb{T}}
%\newcommand{\C}{\mathbb{C}}
\newcommand{\Cone}{\mathbb{C}_{\leq 1}}
\newcommand{\Scheds}{\sfsymbol{Scheds}}


\newcommand{\Pairs}{\ensuremath{\mathbb{P}}}
\newcommand{\EqPairs}{\ensuremath{\mathbb{I\hspace{-.15ex}E}}}

\newcommand{\iwpair}[2]{({#1} \leftsliceeq {#2})}
\newcommand{\EqPair}[2]{\Lbag{#1} \leftsliceeq {#2}\Rbag} 
\newcommand{\littlebigEqPair}[2]{\littlebigLbag{#1} ~{}\mathrel{\text{\large$\leftsliceeq$}}{}~ {#2}\littlebigRbag}
\newcommand{\bigEqPair}[2]{\bigLbag{#1} ~{}\mathrel{\text{\large$\leftsliceeq$}}{}~ {#2}\bigRbag}
\newcommand{\BigEqPair}[2]{\BigLbag{#1} ~{}\mathrel{\text{\large$\leftsliceeq$}}{}~ {#2}\BigRbag}
\newcommand{\hugeEqPair}[2]{\hugeLbag{#1} ~{}\mathrel{\text{\Large$\leftsliceeq$}}{}~ {#2}\hugeRbag}



\newcommand{\DExprs}{\ensuremath{\mathsf{PGuards}}\xspace}    % Distribution Expression
\newcommand{\dists}[1]{\ensuremath{\mathcal{D}\left({#1}\right)}}
\newcommand{\Valdists}{\dists{\Vals}}





%%%%    Functions    %%%%

\newcommand{\dom}[1]{\sfsymbol{dom}\left({#1}\right)}
\newcommand{\iverson}[1]{\left[ {#1} \right]}
\newcommand{\iiverson}[1]{\ensuremath{\Lbag #1 \Rbag}}
\newcommand{\ket}[1]{\left|\, {#1} \right\rangle}
\newcommand{\bigiverson}[1]{\vphantom{\big(}\iverson{#1}\vphantom{\big)}}
\newcommand{\Min}[2]{\min\left\{\,{#1},\: {#2}\,\right\}}
\newcommand{\cemphMin}[2]{\cemph{\min}\left\{\,{#1},\: {#2}\,\right\}}
\newcommand{\Max}[2]{\max\left\{\,{#1},\: {#2}\,\right\}}
%\newcommand{\Exp}[2]{\sfsymbol{Exp}_{#1}\left( {#2} \right)}
\newcommand{\Exp}[2]{\int_\States \; {#2} \; d \,{#1} }
\newcommand{\ExpVals}[2]{\int_\Vals \; {#2} \; d \,{#1} }
\newcommand{\ExpNoSet}[2]{\int \; {#2} \; d \,{#1} }
\newcommand{\ERT}[2]{\sfsymbol{ERT}\left\llbracket #1 \right\rrbracket_{#2}}
\newcommand{\positive}[1]{\tensor*[^{+}]{#1}{}}
\newcommand{\negative}[1]{\tensor*[^{-}]{#1}{}}

\newcommand{\pos}[1]{(#1)^{+}}

%\newcommand{\Exp}[2]{\int_\States {#2} \: d {#1} }

\newcommand{\subst}[2]{\left[ {#1} \middle/ {#2}\right]}
\newcommand{\statesubst}[2]{\left[ {#1} \mapsto {#2}\right]}


\newcommand{\monus}{% Don't redefine it if available
	\mathbin{% We want a binary operation
		\vphantom{+}% The same height as a plus or minus
		\text{% Change size in sub/superscripts
			\mathsurround=0pt % To be on the safe side
			\ooalign{% Superimpose the two symbols
				\noalign{\kern-.5ex}% but the dot is raised a bit
				\hidewidth$\smash{\cdot}$\hidewidth\cr % Dot
				\noalign{\kern.5ex}% Backup for vertical alignment
				$-$\cr % Minus
			}%
		}%
	}%
}

%\newcommand{\charfun}[4]{\tensor*[^{#2}_{{#1}{#3}}]{\Phi}{_{{#4}}}}
%\newcommand{\charfun}[4]{\tensor*[^{#1}_{{\langle #2, #3 \rangle}}]{\Phi}{_{{#4}}}}
%\newcommand{\charfun}[4]{\charfunn{#1}{#2}{#3}{#4}{}}
\newcommand{\charfun}[1]{\Phi_{#1}}
\newcommand{\aertcharfun}[1]{\Psi_{#1}}
%\newcommand{\charfunn}[5]{\tensor*[^{#1}_{\ifthenelse{\equal{#2}{}\and\equal{#3}{}}{}{\langle #2, #3 \rangle}}]{\Phi}{_{{#4}}^{{#5}}}}
%
%\newcommand{\charfunnoguard}[2]{\tensor*[^{#1}]{\Phi}{_{{#2}}}}
%\newcommand{\charfunnoguardn}[3]{\tensor*[^{#1}]{\Phi}{_{{#2}}^{{#3}}}}

\newcommand{\charwp}[3]{\charfun{\wpsymbol}{#1}{#2}{#3}}
\newcommand{\charwlp}[3]{\charfun{\wlpsymbol}{#1}{#2}{#3}}
\newcommand{\charawp}[3]{\charfun{\awpsymbol}{#1}{#2}{#3}}
\newcommand{\charawlp}[3]{\charfun{\awlpsymbol}{#1}{#2}{#3}}
\newcommand{\charert}[3]{\charfun{\ertsymbol}{#1}{#2}{#3}}
\newcommand{\chariwp}[3]{\charfun{\iwpsymbol}{#1}{#2}{#3}}
\newcommand{\charcwp}[3]{\charfun{\cwpsymbol}{#1}{#2}{#3}}
\newcommand{\charcwlp}[3]{\charfun{\cwlpsymbol}{#1}{#2}{#3}}

\newcommand{\charwpn}[4]{\charfunn{\wpsymbol}{#1}{#2}{#3}{#4}}
\newcommand{\charwlpn}[4]{\charfunn{\wlpsymbol}{#1}{#2}{#3}{#4}}
\newcommand{\charawpn}[4]{\charfunn{\awpsymbol}{#1}{#2}{#3}{#4}}
\newcommand{\charawlpn}[4]{\charfunn{\awlpsymbol}{#1}{#2}{#3}{#4}}
\newcommand{\chariwpn}[4]{\charfunn{\iwpsymbol}{#1}{#2}{#3}{#4}}
\newcommand{\charertn}[4]{\charfunn{\ertsymbol}{#1}{#2}{#3}{#4}}
\newcommand{\charcwpn}[4]{\charfunn{\cwpsymbol}{#1}{#2}{#3}{#4}}
\newcommand{\charcwlpn}[4]{\charfunn{\cwlpsymbol}{#1}{#2}{#3}{#4}}



%\newcommand{\eval}[1]{\ensuremath{\llbracket {#1} \rrbracket}}
\newcommand{\probof}[2]{\ensuremath{\llbracket {#1} \colon {#2} \rrbracket}}
\newcommand{\pot}[1]{\mathcal{P}\left({#1}\right)}

\newcommand{\sem}[3]{\left\llbracket #1 \right\rrbracket_{#2}^{#3}}


\newcommand{\varsf}[1]{\sfsymbol{Vars}\left( #1 \right)}
\newcommand{\modC}[1]{\sfsymbol{Mod}\left( #1 \right)}
\newcommand{\notmod}{\mathrel{\not\mathrel{\Cap}}}
\newcommand{\nnotmod}{~{}\notmod{}~}






%%%% Standards %%%%

%%%% Operational %%%%
\newcommand{\ms}{\sigma}
\newcommand{\ma}{a}
\newcommand{\msched}{\mathfrak{S}}

\newcommand{\MDP}{\mathcal{M}}
\newcommand{\MS}{\mathcal{S}}
\newcommand{\MA}{\textsf{Act}}
\newcommand{\MP}{\textsf{prob}}
\newcommand{\MI}{\ms_{\text{init}}}
\renewcommand{\MR}{\textsf{rew}}


\newcommand{\MDPInit}[1]{\MDP[#1]}
\newcommand{\MSCHED}{\textsf{Sched}}

\newcommand{\mstep}[2]{\xrightarrow[#1]{#2}}
\newcommand{\qmstep}[2]{\quad\mstep{#1}{#2}\quad}

\newcommand{\conf}[3]{#1,\, #2,\, #3}
\newcommand{\term}{\mathpzc{term}}
\newcommand{\sink}{\mathpzc{sink}}
\newcommand{\fail}{\mathpzc{fault}}
\newcommand{\tconf}[2]{\conf{\term}{#1}{#2}}

\newcommand{\MPOP}{\mathpzc{prob}}
\newcommand{\MROP}{\mathpzc{rew}}
\newcommand{\MDPOP}[1]{\mathpzc{M}\llbracket #1 \rrbracket}

\newcommand{\mpath}{\pi}
\newcommand{\mpaths}[2]{\textsf{Paths}_{=#1}(#2)}

\newcommand{\rt}{\ensuremath{{\textnormal{\textsf{rt}}}}}
\newcommand{\rtt}[2]{\ensuremath{\rt\llbracket #1 \rrbracket(#2)}}
\newcommand{\oprt}{\ensuremath{{\textnormal{\textsf{oprt}}}}}
\newcommand{\oprtt}[2]{\ensuremath{{\oprt\llbracket #1 \rrbracket(#2)}}}
\newcommand{\eert}{\ensuremath{{\textnormal{\textsf{eert}}}}}
\newcommand{\eertt}[2]{\ensuremath{{\eert\llbracket #1 \rrbracket(#2)}}}

%\newcommand{\sched}{\mathfrak{s}}
%\newcommand{\conf}[6]{\left\langle #1,\, #2,\, #3,\, #4,\, #5,\, #6 \right\rangle}
%\newcommand{\mdp}[4]{\left( #1,\, #2,\, #3,\, #4 \right)}
%\newcommand{\MC}[3]{\left( #1,\, #2,\, #3 \right)}
%\newcommand{\mdpstate}[2]{\left\langle #1,\, #2 \right\rangle}
%\newcommand{\opmdp}[2]{\mathcal{M}^{#1, #2}}
%\newcommand{\opmc}[3]{\mathcal{M}^{#1, #2}_{#3}}
%\newcommand{\mc}[3]{\left\langle #1,\, #2,\, #3 \right\rangle}
%\newcommand{\comptree}[2]{\mathcal{T}^{#1, #2}}
%\newcommand{\sink}[2]{\mathpzc{success}, #1, #2}

%\newcommand{\undesired}{\dagger}

\newcommand{\mdps}{\sigma}
\newcommand{\mdpinit}{\mdps_{\textsf{init}}}

\newcommand{\mdpsched}{\mathfrak{S}}
\newcommand{\mdppath}{\pi}
\newcommand{\mdpa}{a}
\newcommand{\mdpactions}{\textsf{Act}}

\newcommand{\mdpprob}[1]{\MDPP\left(#1\right)}


%\newcommand{\MDPI}[1]{\MDP_{#1}}
\newcommand{\ExpRew}[1]{\textsf{ExpRew}\left(#1\right)}


\newcommand{\ooprel}[2]{~\xrightarrow[#2]{#1}~}
\newcommand{\oprel}[1]{~\xrightarrow{#1}~}
\newcommand{\opdrel}{\oprel{}}
\newcommand{\opnrel}[1]{~\xrightarrow[#1]{}~}
\newcommand{\qoprel}[1]{\quad{\oprel{#1}}\quad}
\newcommand{\qopnrel}[1]{\quad{\opnrel{#1}}\quad}
\newcommand{\qopdrel}{\quad{\opdrel}\quad}

\newcommand{\reward}[1]{\MDPR(#1)}





%%%%    Shortcuts    %%%%

\newcommand{\iffq}{\lqiff}
\newcommand{\iffqq}{\lqqiff}
\newcommand{\impliesq}{\lqimplies}
\newcommand{\impliesqq}{\lqqimplies}

\newcommand{\impliedbyq}{\textnormal{implied by}\quad}
\newcommand{\impliedbyqq}{\textnormal{implied by}\qquad}


\newcommand{\cppreceq}{\blacktrianglelefteq}
\newcommand{\ccppreceq}{~{}~{}\cppreceq{}~{}~}

\newcounter{blacktrianglelefteq}
\newcommand{\blacktrianglelefteq}{%
\mathrel{\stepcounter{blacktrianglelefteq}%
\begin{tikzpicture}%
\node[inner sep=0pt,outer ysep=0.125em, draw=none] (\theblacktrianglelefteq) {$\blacktriangleleft$};
\draw[use as bounding box, draw=none] ($(\theblacktrianglelefteq.north west) + (0em, -.125em)$) rectangle ($(\theblacktrianglelefteq.south east) + (0, 0.0625em)$);
\draw[] ($(\theblacktrianglelefteq.south west) + (.075em, 0em)$) -- ($(\theblacktrianglelefteq.south east) + (-0.075em, 0em)$);
\end{tikzpicture}%
}%
}
\newcommand{\ctpreceq}{\trianglelefteq}
\newcommand{\cctpreceq}{~{}~{}\ctpreceq{}~{}~}

\newcounter{leftsliceeq}
\newcommand{\leftsliceeq}{%
\mathrel{\stepcounter{leftsliceeq}%
\begin{tikzpicture}%
\node[inner sep=0pt,outer sep=0em, draw=none] (\theleftsliceeq) {$\leftslice$};
\draw[use as bounding box, draw=none] ($(\theleftsliceeq.north west) + (0em, -.125em)$) rectangle ($(\theleftsliceeq.south east) + (0, 0.0625em)$);
\draw[] ($(\theleftsliceeq.south west) + (.05em, 0.2em)$) -- ($(\theleftsliceeq.south east) + (-0.225em, -0.00em)$);
\end{tikzpicture}%
}%
}

\newcommand{\pprecsim}{~{}\precsim{}~}
\newcommand{\ssuccsim}{~{}\succsim{}~}
\newcommand{\ssimeq}{~{}\simeq{}~}
\newcommand{\vvee}{~{}\vee{}~}
\newcommand{\wwedge}{~{}\wedge{}~}
\newcommand{\cvee}{\,\vee\,}
\newcommand{\cwedge}{\,\wedge\,}
\newcommand{\qvee}{\quad{}\vee{}\quad}
\newcommand{\qwedge}{\quad{}\wedge{}\quad}
\newcommand{\qqvee}{\qquad{}\vee{}\qquad}
\newcommand{\qqwedge}{\qquad{}\wedge{}\qquad}
\newcommand{\nneq}{~{}\neq{}~}
\newcommand{\aapprox}{~{}\approx{}~}
\newcommand{\pprecapprox}{~\:{}\precapprox{}\:~}
\newcommand{\ssqsupseteq}{~{}\sqsupseteq{}~}
\newcommand{\ssubseteq}{~{}\subseteq{}~}
\newcommand{\ssubset}{~{}\subset{}~}
\newcommand{\ssetminus}{~{}\setminus{}~}
\newcommand{\ccup}{~{}\cup{}~}
\newcommand{\ccap}{~{}\cap{}~}
\newcommand{\pplus}{~{}+{}~}
\newcommand{\rtmin}{\ensuremath{\sqcap}}
\newcommand{\rrtmin}{~{}\ensuremath{\sqcap}~{}}
\newcommand{\mminus}{~{}-{}~}
\newcommand{\ooplus}{~{}\oplus{}~}
\newcommand{\oodot}{~{}\odot{}~}
\newcommand{\qmid}{\quad{}|{}\quad}


\newcommand{\rta}{\ensuremath{f}}
\newcommand{\rtb}{\ensuremath{g}}
\newcommand{\rtc}{\ensuremath{u}}
\newcommand{\inv}{\ensuremath{I}}

\newcommand{\arta}{\ensuremath{X}}
\newcommand{\artb}{\ensuremath{Y}}
\newcommand{\artc}{\ensuremath{Z}}
\newcommand{\ainv}{\ensuremath{J}}

\newcommand{\preda}{\ensuremath{\varphi}} % predicates
\newcommand{\predb}{\ensuremath{\psi}} 
\newcommand{\predc}{\ensuremath{\vartheta}} 
\newcommand{\predd}{\ensuremath{\xi}} 


\newcommand{\pp}{\ensuremath{p}} % probabilities
\newcommand{\qq}{\ensuremath{q}} % probabilities


\newcommand{\nicepar}[1]{\vcenter{\hbox{$\displaystyle #1$}}}

\newcommand{\bigtag}[1]{\tag*{$\left( \text{#1} \right)$}}

\newcommand{\suppreceq}{{\textstyle\sup_{\preceq}}}
\newcommand{\infpreceq}{{\textstyle\inf_{\preceq}}}
\newcommand{\supctpreceq}{{\textstyle\sup_{\ctpreceq}}}
\newcommand{\supcppreceq}{{\textstyle\sup_{\scalebox{0.75}{$\cppreceq$}}}}





\newcommand{\setcomp}[2]{\left\{\, {#1} ~\middle|~ {#2} \,\right\}}

\newcommand{\doublehrule}{\hrule\vspace{2pt}\hrule}
%\newcommand{\doublehrule}{\hrulefill\vspace{2pt}\hrulefill}

\definecolor{webgreen}{rgb}{0,.5,0}

\newcommand{\black}[1]{\textcolor{black}{#1}}
\newcommand{\gray}[1]{\textcolor{gray}{#1}}
\newcommand{\lgray}[1]{\textcolor{lightgray}{#1}}
\newcommand{\green}[1]{\textcolor{webgreen}{#1}}
\newcommand{\red}[1]{\textcolor{Maroon}{#1}}
%\colorlet{myblue}{DodgerBlue3}
\newcommand{\blue}[1]{\textcolor{myblue}{#1}}
\newcommand{\orange}[1]{\textcolor{DarkOrange3}{#1}}
\newcommand{\brown}[1]{\textcolor{webbrown}{#1}}
\newcommand{\darkorchid}[1]{\textcolor{HotPink4}{#1}}
%\newcommand{\annocolor}[1]{\textcolor{DarkOrange3}{#1}}
%\newcommand{\annocolor}[1]{\cemph{#1}}
\newcommand{\definize}[1]{\blue{#1}}
\newcommand{\theoremize}[1]{{#1}}
%\colorlet{cemphcolor}{HotPink4}
%\definecolor{mydarkorange}{RGB}{253,97,0}
\definecolor{mydarkorange}{RGB}{220,80,0}
\colorlet{cemphcolor}{mydarkorange}
%\colorlet{cemphcolor}{Maroon}
\newcommand{\cemph}[1]{\textcolor{cemphcolor}{#1}}

\newcommand{\underdot}[1]{%
    \tikz[baseline=(todotted.base)]{
\node[inner sep=1pt,outer sep=0pt] (todotted) {#1};
\draw[densely dotted] (todotted.south west) -- (todotted.south east);
    }%
}%

\newcommand{\cloze}[1]{\underdot{\hphantom{$#1$}}}


\newcommand{\ruletext}[1]{\textsf{#1}}

\newcommand{\annocolor}[1]{\textcolor{DodgerBlue3}{#1}}
\newcommand{\aannocolor}[1]{\textcolor{webgreen}{#1}}
%\newcommand{\annocolor}[1]{\textcolor{blue}{#1}}
\newcommand{\annotate}[1]{{\annocolor{\!\!{\fatslash}\!\!{\fatslash}~~\vphantom{G'} {#1}}}}
\newcommand{\sndannotate}[1]{{\annocolor{\vphantom{G'} \qquad{#1}}}}
\newcommand{\starannotate}[1]{{\annocolor{{\talloblong}\!{\talloblong}\:\vphantom{G'} {#1}}}}
\newcommand{\aertstarannotate}[1]{{\aannocolor{\!\!\hspace{-0.0ex}{}^{\aannocolor{{\text{\tiny $\potent$}}}}{\talloblong}\!{\talloblong}\:\vphantom{G'} {#1}}}}
\newcommand{\phiannotate}[1]{{\annocolor{\!\!\hspace{-.15ex}{}^{\annocolor{\Phi}}\!\!\!{\fatslash}\!\!{\fatslash}~~\vphantom{G'} {#1}}}}

\newcommand{\psiannotate}[1]{{\aannocolor{\!\!\hspace{-.15ex}{}^{\aannocolor{\Psi}}\!\!\!{\fatslash}\!\!{\fatslash}~~\vphantom{G'} {#1}}}}

\newcommand{\bowtieannotate}[1]{\annocolor{\!\!\hspace{-.35ex}{}^{\annocolor{{\bowtie}}}{\!\!\!{\fatslash}\!\!{\fatslash}~~\vphantom{G'} {#1}}}}

\newcommand{\preceqannotate}[1]{\annocolor{\!\!\hspace{-.25ex}{}^{\annocolor{{\preceq}}}{\!\!\!{\fatslash}\!\!{\fatslash}~~\vphantom{G'} {#1}}}}

\newcommand{\geqannotate}[1]{\annocolor{\!\!\hspace{-.1ex}{}^{\annocolor{\text{\rotatebox{-90}{$\geq$}}}}{\!\!\!{\fatslash}\!\!{\fatslash}~~\vphantom{G'} {#1}}}}

\newcommand{\mutannotate}[1]{\annocolor{\!\!\hspace{-1.3ex}{}^{\annocolor{{\text{\tiny $\ruletext{mut}$}}}}{\!\!\!{\fatslash}\!\!{\fatslash}~~\vphantom{G'} {#1}}}}

\newcommand{\lkpannotate}[1]{\annocolor{\!\!\hspace{-1.3ex}{}^{\annocolor{{\text{\tiny $\ruletext{lkp}$}}}}{\!\!\!{\fatslash}\!\!{\fatslash}~~\vphantom{G'} {#1}}}}

\newcommand{\succeqannotate}[1]{\annocolor{\!\!\hspace{-.25ex}{}^{\annocolor{{\succeq}}}{\!\!\!{\fatslash}\!\!{\fatslash}~~\vphantom{G'} {#1}}}}

\newcommand{\eqannotate}[1]{\annocolor{\!\!\hspace{-.1ex}{}^{\annocolor{{=}}}{\!\!\!{\fatslash}\!\!{\fatslash}~~\vphantom{G'} {#1}}}}

\newcommand{\frameannotate}[2]{{\annocolor{\!\!\hspace{-2.25ex}{}^{\annocolor{\text{\tiny $\ruletext{frame}$}}}_{#2}\!\!\!{\fatslash}\!\!{\fatslash}~~\vphantom{G'} {#1}}}}

\newcommand{\unframeannotate}[2]{{\annocolor{\!\!\hspace{-4.0ex}{}^{\annocolor{\text{\tiny $\ruletext{unframe}$}}}_{#2}\!\!\!{\fatslash}\!\!{\fatslash}~~\vphantom{G'} {#1}}}}

\newcommand{\framelkpannotate}[1]{{\annocolor{\!\!\hspace{-2.25ex}{}^{\annocolor{\text{\tiny $\ruletext{frame}$}}}_{\text{\tiny $\ruletext{lkp}$}}\!\!\!{\fatslash}\!\!{\fatslash}~~\vphantom{G'} {#1}}}}



\newcommand{\framemutannotate}[1]{{\annocolor{\!\!\hspace{-2.25ex}{}^{\annocolor{\text{\tiny $\ruletext{frame}$}}}_{\text{\tiny $\ruletext{mut}$}}\!\!\!{\fatslash}\!\!{\fatslash}~~\vphantom{G'} {#1}}}}

\newcommand{\auxvarannotate}[1]{{\annocolor{\!\!\hspace{-0.90ex}{}^{\annocolor{\text{\tiny $\ruletext{aux}$}}}\!\!\!{\fatslash}\!\!{\fatslash}~~\vphantom{G'} {#1}}}}

\newcommand{\ldotannotate}[1]{{\annocolor{\!\!\hspace{-0.5ex}{}^{\annocolor{\text{\tiny $\ldots$}}}\!\!\!{\fatslash}\!\!{\fatslash}~~\vphantom{G'} {#1}}}}

\newcommand{\ertannotate}[1]{{\annocolor{\!\!\hspace{-0.5ex}{}^{\annocolor{\text{\tiny $\ertsymbol$}}}\!\!\!{\fatslash}\!\!{\fatslash}~~\vphantom{G'} {#1}}}}


\newcommand{\aertannotate}[1]{{\aannocolor{\!\!\hspace{0.1ex}{}^{\aannocolor{\text{\tiny $\potent$}}}\!\!\!{\fatslash}\!\!{\fatslash}~~\vphantom{G'} {#1}}}}

\newcommand{\wpannotate}[1]{{\annocolor{\!\!\hspace{-0.75ex}{}^{\annocolor{\text{\tiny $\wpsymbol$}}}\!\!\!{\fatslash}\!\!{\fatslash}~~\vphantom{G'} {#1}}}}

\newcommand{\wlpannotate}[1]{{\annocolor{\!\!\hspace{-0.75ex}{}^{\annocolor{\text{\tiny $\wlpsymbol$}}}\!\!\!{\fatslash}\!\!{\fatslash}~~\vphantom{G'} {#1}}}}

\newcommand{\phantannotate}[1]{{\annocolor{\vphantom{\!\!{\fatslash}\!\!{\fatslash}~~}\vphantom{G'} {#1}}}}





\newcommand{\aaertannotate}[1]{{\aannocolor{\!\!\hspace{-1.3ex}{}^{\aannocolor{\text{\tiny $\aertsymbol$}}}\!\!\!{\fatslash}\!\!{\fatslash}~~\vphantom{G'} {#1}}}}

\newcommand{\aannotate}[1]{{\aannocolor{\!\!{\fatslash}\!\!{\fatslash}~~\vphantom{G'} {#1}}}}

\newcommand{\asucceqannotate}[1]{\aannocolor{\!\!\hspace{-.25ex}{}^{\aannocolor{{\succeq}}}{\!\!\!{\fatslash}\!\!{\fatslash}~~\vphantom{G'} {#1}}}}


\newcommand{\aframeannotate}[1]{{\aannocolor{\!\!\hspace{-2.25ex}{}^{\aannocolor{\text{\tiny $\ruletext{frame}$}}}\!\!\!{\fatslash}\!\!{\fatslash}~~\vphantom{G'} {#1}}}}

\newcommand{\aunframeannotate}[1]{{\aannocolor{\!\!\hspace{-4.0ex}{}^{\aannocolor{\text{\tiny $\ruletext{unframe}$}}}\!\!\!{\fatslash}\!\!{\fatslash}~~\vphantom{G'} {#1}}}}


\newcommand{\amutannotate}[1]{\aannocolor{\!\!\hspace{-1.3ex}{}^{\aannocolor{{\text{\tiny $\ruletext{mut}$}}}}{\!\!\!{\fatslash}\!\!{\fatslash}~~\vphantom{G'} {#1}}}}

\newcommand{\astarannotate}[1]{{\aannocolor{{\talloblong}\!{\talloblong}\:\vphantom{G'} {#1}}}}


\newcommand{\aphantannotate}[1]{{\aannocolor{\vphantom{\!\!{\fatslash}\!\!{\fatslash}~~}\vphantom{G'} {#1}}}}


\newcommand{\alkpannotate}[1]{\aannocolor{\!\!\hspace{-1.3ex}{}^{\aannocolor{{\text{\tiny $\ruletext{lkp}$}}}}{\!\!\!{\fatslash}\!\!{\fatslash}~~\vphantom{G'} {#1}}}}


\newcommand{\aauxvarannotate}[1]{{\aannocolor{\!\!\hspace{-0.90ex}{}^{\aannocolor{\text{\tiny $\ruletext{aux}$}}}\!\!\!{\fatslash}\!\!{\fatslash}~~\vphantom{G'} {#1}}}}


% helpers
\newcommand{\phwhere}{\phantom{\text{where }}}
\newcommand{\phand}{\phantom{\text{and }}}

\newcommand{\exprog}[1]{\ensuremath{C_\textnormal{\ref*{#1}}}}
\newcommand{\subexprog}[2]{\ensuremath{C_\textnormal{\ref*{#1}~\ref*{#2}.}}}

\newcounter{computationarrowsone}
\newcounter{computationarrowstwo}
\newcommand{\tikzcomputationarrow}[4]{%
	\stepcounter{computationarrowsone}
	\stepcounter{computationarrowstwo}
	\draw [decorate,-{Stealth[scale=1]}] (#1,#2) -- (#3,#4);%
	\node (\thecomputationarrowsone) at (#1,#2) {};%
	\fill (\thecomputationarrowsone) + (-.06,0) circle [radius=1.75pt];%
	\node (\thecomputationarrowstwo) at (#3,#4) {};%
	\fill (\thecomputationarrowstwo) + (.06,0) circle [radius=1.75pt];}
	
\newcommand{\tikzbendleftcomputationarrow}[4]{%
	\stepcounter{computationarrowsone}
	\stepcounter{computationarrowstwo}
	\draw (#1,#2) edge[decorate,-{Stealth[scale=1]},bend left] (#3,#4);%
	\node (\thecomputationarrowsone) at (#1,#2) {};%
	\fill (\thecomputationarrowsone) + (-.06,0) circle [radius=1.75pt];%
	\node (\thecomputationarrowstwo) at (#3,#4) {};%
	\fill (\thecomputationarrowstwo) + (.06,0) circle [radius=1.75pt];}
	
\newcommand{\tikzbendrightcomputationarrow}[4]{%
	\stepcounter{computationarrowsone}
	\stepcounter{computationarrowstwo}
	\draw (#1,#2) edge[decorate,-{Stealth[scale=1]},bend right] (#3,#4);%
	\node (\thecomputationarrowsone) at (#1,#2) {};%
	\fill (\thecomputationarrowsone) + (-.06,0) circle [radius=1.75pt];%
	\node (\thecomputationarrowstwo) at (#3,#4) {};%
	\fill (\thecomputationarrowstwo) + (.06,0) circle [radius=1.75pt];}
	
\newcommand{\tikzpossiblecomputationarrow}[4]{%
	\stepcounter{computationarrowsone}
	\stepcounter{computationarrowstwo}
	\draw [decorate,-{Stealth[scale=1]},dashdotted] (#1,#2) -- (#3,#4);%
	\node (\thecomputationarrowsone) at (#1,#2) {};%
	\fill (\thecomputationarrowsone) + (-.06,0) circle [radius=1.75pt];%
	\node (\thecomputationarrowstwo) at (#3,#4) {};%
	\fill (\thecomputationarrowstwo) + (.06,0) circle [radius=1.75pt];}

%\newcommand{\drule}[3]{\left( \text{#1} \right)~~\frac{~~{#2}~~}{~~{#3}~~}}
\newcommand{\drule}[3]{\infer[\left( \textnormal{#1} \right)]{~~~~{#3}~~~~}{~~~~{#2}~~~~}}
\newcommand{\drulecomma}[3]{\infer[\left( \textnormal{#1} \right)~,]{~~~~{#3}~~~~}{~~~~{#2}~~~~}}
\newcommand{\druledot}[3]{\infer[\left( \textnormal{#1} \right)~.]{~~~~{#3}~~~~}{~~~~{#2}~~~~}}
\newcommand{\drulenolabel}[2]{\infer{~~{#2}~~}{~~{#1}~~}}


\newcommand{\confto}{\vdash}
\newcommand{\cconfto}{~{}\confto{}~}
\newcommand{\conftor}{\mathrel{{\confto}^*}}
\newcommand{\cconftor}{~{}\conftor{}~}
\newcommand{\conftok}[1]{\mathrel{{\confto}^{#1}}}
\newcommand{\cconftok}[1]{~{}\conftok{#1}{}~}
\newcommand{\conftosched}[1]{\confto_{#1}}
\newcommand{\cconftosched}[1]{~{}\conftosched{#1}{}~}
\newcommand{\conftoschedr}[1]{\mathrel{\confto_{#1}^*}}
\newcommand{\cconftoschedr}[1]{~{}\conftoschedr{#1}{}~}
\newcommand{\xconfto}[1]{\stackrel{#1}{\confto}}
\newcommand{\xxconfto}[1]{~{}\xconfto{#1}{}~}

\newcommand{\phantoma}{\phantom{\textsf{a}}}
\newcommand{\phantoml}{\phantom{\textsf{l}}}
\newcommand{\phantomal}{\phantom{\textsf{al}}}

\newcommand{\pautoref}[1]{\textnormal{\autoref{#1}} on p.~\pageref{#1}}
\newcommand{\pthref}[1]{\textnormal{\thref{#1}} on p.~\pageref{#1}}
\newcommand{\psubthref}[2]{\textnormal{\subthref{#1}{#2}} on p.~\pageref{#1:#2}}







%%%%    Symbols    %%%%

\newcounter{sarrow}
\newcommand\xrsquigarrow[1]{%
\stepcounter{sarrow}%
\begin{tikzpicture}[decoration=snake]
\node (\thesarrow) {\strut{#1}};
\draw[->, decorate] (\thesarrow.south west) -- (\thesarrow.south east);
\end{tikzpicture}%
}

\newcommand{\lfp}{\ensuremath{\textnormal{\sfsymbol{lfp}}~}}
\newcommand{\gfp}{\ensuremath{\textnormal{\sfsymbol{gfp}}~}}






%%%%% TERRIBLE HACKS %%%%%

\newcommand{\lowerpenalties}{\allowdisplaybreaks\predisplaypenalty=0\clubpenalty=0\widowpenalty=0\displaywidowpenalty=0\brokenpenalty=0{}}

\newcommand{\restorepenalties}{\allowdisplaybreaks[0]\predisplaypenalty=10000\clubpenalty=150\widowpenalty=150\displaywidowpenalty=50\brokenpenalty=100{}}

\newcommand{\unlowerpenalties}{\restorepenalties}





%%%%% TERRIBLE HACKS FROM THE INTERNET %%%%%

\makeatletter
\newlength{\negph@wd}
\DeclareRobustCommand{\negphantom}[1]{%
  \ifmmode
    \mathpalette\negph@math{#1}%
  \else
    \negph@do{#1}%
  \fi
}
\newcommand{\negph@math}[2]{\negph@do{$\m@th#1#2$}}
\newcommand{\negph@do}[1]{%
  \settowidth{\negph@wd}{#1}%
  \hspace*{-\negph@wd}%
}
\makeatother




% From the ACTA paper


\newcommand{\Configs}{\mathbb{K}}


\newcommand{\kleeneT}{\ensuremath{\textnormal T}\xspace}
\newcommand{\Tp}{\ensuremath{\T_p}\xspace}
\newcommand{\Tps}{\ensuremath{\Tp^*}\xspace}
\newcommand{\ExpN}[3]{\ensuremath{\textnormal E_{#1}^{#3}\left({#2}\right)}\xspace}
\renewcommand{\Pr}{\textnormal{Pr}}
\newcommand{\covsymbol}{\mathnormal{\mathsf{Cov}}}
\newcommand{\Cov}[4]{\covsymbol_{\llbracket #1 \rrbracket (#2) }\left( {#3},\, {#4} \right)}
\newcommand{\cCov}[3]{\covsymbol\left( {#1},\, {#2} \mathrel{|} {#3} \right)}

\newcommand{\toppr}{\ensuremath{\bot_{\textnormal{pr}}}}
\newcommand{\topnd}{\ensuremath{\bot_{\textnormal{nd}}}}
\newcommand{\topter}{\ensuremath{\bot_{\textnormal{ter}}}}

\newcommand{\Problem}[1]{\mathcal{#1}}
\newcommand{\cProblem}[1]{\overline{\text{\footnotesize{$\Problem{#1}$}}}}

\newcommand{\HP}{\Problem{H}\xspace}
\newcommand{\cHP}{\cProblem{H}\xspace}
\newcommand{\UHP}{\,\Problem{UH}\xspace}
\newcommand{\cUHP}{\,\cProblem{UH}\xspace}
\newcommand{\COF}{\Problem{COF}\xspace}
\newcommand{\cCOF}{\cProblem{COF}\xspace}
\newcommand{\AST}{\ensuremath{\Problem{AST}}\xspace}
\newcommand{\cAST}{\ensuremath{\cProblem{AST}}\xspace}
\newcommand{\UAST}{\ensuremath{\Problem{U\hspace{-1pt}AST}}\xspace}
\newcommand{\cUAST}{\ensuremath{\cProblem{U\hspace{-1pt}AST}}\xspace}
\newcommand{\PAST}{\ensuremath{\Problem{P\hspace{-2pt}AST}}\xspace}
\newcommand{\cPAST}{\ensuremath{\cProblem{P\hspace{-2pt}AST}}\xspace}
\newcommand{\UPAST}{\ensuremath{\Problem{UP\hspace{-2pt}AST}}\xspace}
\newcommand{\cUPAST}{\ensuremath{\cProblem{UP\hspace{-2pt}AST}}\xspace}
\newcommand{\EXP}{\ensuremath{\Problem{E\hspace{-1pt}XP}}}
\newcommand{\LEXP}{\ensuremath{\Problem{LE\hspace{-1pt}XP}}}
\newcommand{\REXP}{\ensuremath{\Problem{RE\hspace{-1pt}XP}}}
\newcommand{\FEXP}{\ensuremath{\Problem{F\hspace{-.5pt}E\hspace{-1pt}XP}}}
\newcommand{\DCOVAR}{\ensuremath{\Problem{DCOV\hspace{-.5ex}AR}}}
\newcommand{\LCOVAR}{\ensuremath{\Problem{LCOV\hspace{-.5ex}AR}}}
\newcommand{\RCOVAR}{\ensuremath{\Problem{RCOV\hspace{-.5ex}AR}}}
\newcommand{\LVAR}{\ensuremath{\Problem{LV\hspace{-.5ex}AR}}}
\newcommand{\RVAR}{\ensuremath{\Problem{RV\hspace{-.5ex}AR}}}
\newcommand{\COVAR}{\ensuremath{\Problem{COV\hspace{-.5ex}AR}}}
\newcommand{\VAR}{\ensuremath{\Problem{V\hspace{-.5ex}AR}}}
\newcommand{\INFCOVAR}{\ensuremath{{}^{\infty}\Problem{COV\hspace{-.5ex}AR}}}
\newcommand{\INFVAR}{\ensuremath{{}^\infty\Problem{V\hspace{-.5ex}AR}}}
\newcommand{\LRTVAR}{\ensuremath{\Problem{L\hspace{-.25ex}R\hspace{-.25ex}TV\hspace{-.5ex}AR}}}
\newcommand{\RRTVAR}{\ensuremath{\Problem{R\hspace{-.25ex}R\hspace{-.25ex}TV\hspace{-.5ex}AR}}}
\newcommand{\RTVAR}{\ensuremath{\Problem{R\hspace{-.25ex}TV\hspace{-.5ex}AR}}}
\newcommand{\INFRTVAR}{\ensuremath{{}^\infty\Problem{R\hspace{-.25ex}TV\hspace{-.5ex}AR}}}

\newcommand{\leqm}{\mathrel{\:{\leq}_{\textnormal{m}}}}
\newcommand{\leqT}{\mathrel{\:{\leq}_{\textnormal{T}}}}
\newcommand{\equivm}{\mathrel{\:{\equiv}_{\textnormal{m}}}}
\newcommand{\equivT}{\mathrel{\:{\equiv}_{\textnormal{T}}}}

\newcommand{\LLL}[1]{\ensuremath{\textnormal{L}\big(#1\big)}\xspace}
\newcommand{\RRR}[1]{\ensuremath{\textnormal{R}\big(#1\big)}\xspace}


% For aert part
\newcommand{\aertsymbol}{\ensuremath{{\textnormal{\textsf{aert}}}}}
% \newcommand{\aert}[4]{\aertsymbol_{#1} \left\llbracket{#2} \right\rrbracket \EqPair{#3}{#4}}
\newcommand{\aert}[3]{\aertOne{#1}{#2}{#3}} % does not use the pairs anymore
\newcommand{\aertOne}[3]{\aertsymbol_{#1}  \left\llbracket{#2} \right\rrbracket \left({#3}\right)}
\newcommand{\aertfunction}[2]{\aertsymbol_{#1}  \left\llbracket{#2}\right\rrbracket}
\newcommand{\boldaertsymbol}{\ensuremath{\boldsymbol{{\textnormal{\textbf{\textsf{aert}}}}}}}

\newcommand{\potent}{\ensuremath{\pi}}

\newcommand{\iwpOne}[2]{\iwpsymbol \left\llbracket{#1} \right\rrbracket \left({#2}\right)}
\newcommand{\iwpfunction}[1]{\iwpsymbol \left\llbracket{#1} \right\rrbracket}
\newcommand{\ertfunction}[1]{\ertsymbol \left\llbracket{#1} \right\rrbracket}
\newcommand{\wpfunction}[1]{\wpsymbol \left\llbracket{#1} \right\rrbracket}


% First and second component
\newcommand{\co}[1]{\ensuremath{\text{first} \left(#1\right)}}
\newcommand{\second}[1]{\ensuremath{\text{second} \left(#1\right)}}

%Least fixpoint
\newcommand{\lfpX}[1]{\ensuremath{\text{lfp}\ X . #1}}

% Program components
\newcommand{\askip}{\ensuremath{{\textnormal{\textsf{skip}}}}}
\newcommand{\diverge}{\ensuremath{{\textnormal{\textsf{diverge}}}}}
\newcommand{\assign}{\ensuremath{x \coloneqq E}}
\newcommand{\tick}{\ensuremath{{\textnormal{\textsf{tick(n)}}}}}
\newcommand{\randassign}{\ensuremath{x \colonapprox \mu}}
\newcommand{\ifelse}{\ensuremath{\ensuremath{{\textnormal{\textsf{if}}}}(\varphi)\{C_1\} \ensuremath{{\textnormal{\textsf{ else}}}} \{C_2\}}}
\newcommand{\prob}{\ensuremath{\{C_1\}[p]\{C_2\}}}
\newcommand{\concat}{\ensuremath{C_1 ; C_2}}
\newcommand{\while}{\ensuremath{\ensuremath{{\textnormal{\textsf{while}}}} (\varphi)\{C'\}}}
\newcommand{\whilepl}[2]{\ensuremath{\ensuremath{{\textnormal{\textsf{while}}}}(#1)\{#2\}}}

\newcommand{\ifC}[1]{\textnormal{\textsf{if}}(\varphi)\{#1\}}


%Characteristic functions
\newcommand{\charaert}[2]{\ensuremath{\prescript{\aertsymbol_{\pi}}{<\varphi,C'>}{\phi_{#1} \left(#2\right)}}}
\newcommand{\charaertno}[2]{\ensuremath{\prescript{\aertsymbol_{\pi}}{<\varphi,C'>}{\phi_{#1} #2}}}
\newcommand{\chariwpno}[2]{\ensuremath{\prescript{\iwpsymbol_{\pi}}{<\varphi,C'>}{\phi_{#1} #2}}}
\newcommand{\chariwpnosh}[2]{\ensuremath{\prescript{\iwpsymbol_{\pi}}{}{\phi_{#1} #2}}}

\newcommand{\charaertn}[3]{\ensuremath{\prescript{\aertsymbol_{\pi}}{<\varphi,C'>}{\phi^{#1}_{#2} #3}}}

% char without <phi,C> for readibility
\newcommand{\charaertnsh}[3]{\ensuremath{\prescript{\aertsymbol_{\pi}}{}{\phi^{#1}_{#2} (#3)}}}
\newcommand{\charaertsh}[2]{\ensuremath{\prescript{\aertsymbol_{\pi}}{}{\phi_{#1}#2}}}
\newcommand{\charwpnsh}[3]{\ensuremath{\prescript{\wpsymbol}{}{\phi^{#1}_{#2} \left(#3\right)}}}
\newcommand{\charertnsh}[3]{\ensuremath{\prescript{\ertsymbol}{}{\phi^{#1}_{#2} \left(#3\right)}}}
\newcommand{\charertsh}[2]{\ensuremath{\prescript{\ertsymbol}{}{\phi_{#1} \left(#2\right)}}}
\newcommand{\chariwpnsh}[3]{\ensuremath{\prescript{\iwpsymbol}{}{\phi^{#1}_{#2} #3}}}

\newcommand{\charaertnColor}[3]{\ensuremath{\prescript{\genColor{\aertsymbol_{\pi}}}{<\varphi,C'>}{\phi^{#1}_{#2} #3}}}
\newcommand{\charwpnColor}[3]{\ensuremath{\prescript{\genColor{\wpsymbol}}{<\varphi,C'>}{\phi^{#1}_{#2} \left(#3\right)}}}
\newcommand{\charertnColor}[3]{\ensuremath{\prescript{\genColor{\ertsymbol}}{<\varphi,C'>}{\phi^{#1}_{#2} \left(#3\right)}}}
\newcommand{\chariwpnColor}[3]{\ensuremath{\prescript{\genColor{\iwpsymbol}}{<\varphi,C'>}{\phi^{#1}_{#2} #3}}}


\newcommand{\charwpfunct}[1]{\ensuremath{\prescript{\wpsymbol}{<\varphi,C'>}{\phi_{#1}}}}
\newcommand{\charertfunct}[1]{\ensuremath{\prescript{\ertsymbol}{<\varphi,C'>}{\phi_{#1}}}}
\newcommand{\charertfunctsh}[1]{\ensuremath{\prescript{\ertsymbol}{}{\phi_{#1}}}}
\newcommand{\charaertfunctsh}[2]{\ensuremath{\prescript{\aertsymbol_{\pi}}{}{\phi_{\EqPair{#1}{#2}}}}}

\newcommand{\randassignRule}[1]{\ensuremath{\lambda \sigma .\text{\, } \int_{Vals} \left( \lambda v.\text{\, }#1\left(\sigma[x \to v ] \right) \right) d \mu_\sigma}}
\newcommand{\randassignshort}[1]{\ensuremath{\vartheta \left(#1\right) }}

\newcommand{\Sup}{\reflectbox{\textnormal{\textsf{\fontfamily{phv}\selectfont S}}}\hspace{.2ex}}
\newcommand{\Inf}{\raisebox{.6\depth}{\rotatebox{-30}{\textnormal{\textsf{\fontfamily{phv}\selectfont \reflectbox{J}}}}\hspace{-.1ex}}}

% for texts, backslash because empty space is deleted
\newcommand{\pGCL}{\text{pGCL}}

\newcommand{\ie}{i.e.\xspace}
\newcommand{\Ie}{I.\,e.\xspace}
\newcommand{\eg}{e.\,g.\xspace}
\newcommand{\Eg}{E.\,g.\xspace} 

\newcommand{\True}{\ensuremath{\mathit{true}}}
\newcommand{\False}{\ensuremath{\mathit{false}}}

\newcommand{\kapit}[1]{\ensuremath{\textnormal{\textsc{#1}}} \ }


% colored EqPairs
\newcommand{\EqPairColor}[2]{\textcolor{red!55!black}{\Lbag{#1} \leftsliceeq {#2}\Rbag}}
\newcommand{\aertColor}[4]{\textcolor{red!55!black}{\aertsymbol_{#1} \left\llbracket{#2} \right\rrbracket \EqPair{#3}{#4}}}
\newcommand{\genColor}[1]{\textcolor{red!55!black}{#1}}
% commands in program listings -> use \annotate{} instead
\newcommand{\comm}[1]{\textcolor{darkgray}{// #1}}

\newcommand{\potex}[1]{\max\{#1,0\}}

% Renamed: were charwpn etc.
\newcommand{\charwpnnew}[3]{\ensuremath{\prescript{\wpsymbol}{<\varphi,C'>}{\phi^{#1}_{#2} \left(#3\right)}}}
\newcommand{\charertnnew}[3]{\ensuremath{\prescript{\ertsymbol}{<\varphi,C'>}{\phi^{#1}_{#2} \left(#3\right)}}}
\newcommand{\chariwpnnew}[3]{\ensuremath{\prescript{\iwpsymbol}{<\varphi,C'>}{\phi^{#1}_{#2} #3}}}

\newcommand{\charwpnew}[2]{\ensuremath{\prescript{\wpsymbol}{<\varphi,C'>}{\phi_{#1} \left(#2\right)}}}
\newcommand{\charertnew}[2]{\ensuremath{\prescript{\ertsymbol}{<\varphi,C'>}{\phi_{#1} \left(#2\right)}}}
\newcommand{\chariwpnew}[2]{\ensuremath{\prescript{\iwpsymbol_{\pi}}{<\varphi,C'>}{\phi_{#1} \left(#2\right)}}}

% Spacing
\newcommand{\qqquad}{\quad \quad \quad}
\newcommand{\qqqquad}{\quad \quad \quad \quad}
\newcommand{\qqqqquad}{\quad \quad \quad \quad \quad}


% Programs
\newcommand{\procfont}[1]{\textit{#1}}

\newcommand{\listtraverse}[1]{C}

\newcommand{\listinsert}[1]{\ensuremath{\procfont{Insert(#1)}}}
\newcommand{\arraycopy}[3]{\ensuremath{\procfont{ArrayCopy(#1,\, #2,\, #3)}}}
\newcommand{\arraydelete}[2]{\ensuremath{\procfont{DeleteArray(#1,\, #2)}}}

\newcommand{\ahead}{A}
\newcommand{\aoff}{o}
\newcommand{\asize}{s}

% Special Expectations
\newcommand{\expfont}[1]{\textsf{#1}}
\newcommand{\elist}[2]{\expfont{list}\left( #1,\, #2 \right)}
\newcommand{\esize}[2]{\expfont{listsize}\left( #1,\, #2 \right)}

\newcommand{\earray}[2]{\expfont{array}\left(#1,\, #2\right)}

\newcommand{\lh}{H}
\newcommand{\ls}{len}
\newcommand{\lpre}{pre}
\newcommand{\lend}{end}
\newcommand{\lsucc}{succ}
%\newcommand{\dllold}[4]{\expfont{dll\textcolor{red}{o}}\left(#1,\, #2, \, #3, \, #4\right)}
\newcommand{\dll}[5]{\expfont{dll}\left(#1,\, #2, \, #3, \, #4, \, #5 \right)}

%FindAny
\newcommand{\procsample}{\procfont{Sample}}
\newcommand{\procrank}{\procfont{Rank}}
\newcommand{\procremove}[1]{\procfont{remove} \left( #1 \right)}
\newcommand{\procdelete}[1]{\procfont{delete} \left( #1 \right)}
\newcommand{\procadd}[1]{\procfont{add} \left( #1 \right)}
\newcommand{\procinsert}[1]{\procfont{insert} \left( #1 \right)}
\newcommand{\procfindany}{\procfont{FindAny}}
\newcommand{\varrank}{rank}
\newcommand{\varany}{any}
\newcommand{\varremove}{x}
\newcommand{\varadd}{y}
\newcommand{\varnew}{N}
\newcommand{\varvalnew}{v_N}
\newcommand{\varvalrem}{v_x}
\newcommand{\varvalany}{v_{any}}
\newcommand{\varvaladd}{v_y}
\newcommand{\varcur}{c}
\newcommand{\varvalcur}{v_c}
\newcommand{\varsamplen}{\ell}
\newcommand{\varlistelem}{y}

\newcommand{\poww}[1]{\textsf{pow2}\left( #1 \right)}
\newcommand{\nextpoww}[1]{\textsf{nextpow2}\left( #1 \right)}

\newcommand{\inputvec}[1]{\overline{#1}}

\newcommand{\proclistinsert}[1]{\procfont{ListInsert}\left( #1\right)}
\newcommand{\procrandlistinsert}[1]{\procfont{RandInsert}\left( #1\right)}


\newcommand{\BOOL}{\ensuremath{\texttt{Bool}}}
\newcommand{\VAL}{\ensuremath{\mathit{Val}}}
\newcommand{\VARS}{\ensuremath{\mathit{Var}}}
\newcommand{\TRUE}{\ensuremath{\texttt{T}}}
\newcommand{\FALSE}{\ensuremath{\texttt{F}}}

\newcommand{\AEXP}{\ensuremath{\mathit{aexp}}}
\newcommand{\FST}{\ensuremath{\texttt{fst}\,}}
\newcommand{\SND}{\ensuremath{\texttt{snd}\,}}
\newcommand{\LET}[2]{\ensuremath{\texttt{let}\; #1 \;\texttt{in}\; #2}}
\newcommand{\FLIP}[1]{\ensuremath{\texttt{flip} \left(#1\right) }}
\newcommand{\NFLIP}{\ensuremath{\texttt{nflip}\left(\right) }}
%\newcommand{\FUNC}[2]{\ensuremath{\texttt{fun}  #1  \left{ {#2} \right} }}
\newcommand{\FUNC}[2]{\ensuremath{\texttt{fun} \; #1 \; \left\{\, #2 \,\right\} } }
\newcommand{\FUNCNAME}{\mathit{func}}




\newcommand{\NDICE}{\sfsymbol{\mbox{noDice}}\xspace}
\newcommand{\DICE}{\sfsymbol{\mbox{Dice}}\xspace}
